\lektion{Welche Entscheidung trifft die Linienmethode?}{molweg}{Methodenwahl}{Vom Feld zum Solver}
\textbf{Auftrag.} Für einen dünnen Stab sollen die Temperaturen in der Mitte und nahe dem rechten Ende als Funktion der Zeit berechnet werden. In der Feldgleichung \(T_t=\alpha T_{xx}\) bezeichnet \(T(x,t)\) noch eine Temperatur an jedem Ort und zu jeder Zeit.
\subt{Was übergeben wir an den Zeitintegrator?}
Wählen wir drei gleich weit entfernte innere Messorte und halten beide Enden bei \(0\,^{\circ}\mathrm C\), so wird aus dem Feld ein Vektor \(U=(T_1,T_2,T_3)^\top\). Die \emph{method of lines} (MOL, Linienmethode) ersetzt die räumlichen Ableitungen. An jedem festen Ort bleibt eine stetige Zeitvariable. Für zentrale Differenzen ergibt sich
\[
 \dot U=\frac{\alpha}{\Delta x^2}
 \begin{pmatrix}-2&1&0\\1&-2&1\\0&1&-2\end{pmatrix}U.
\]
Bei \(U=(0,10,0)^\top\,\mathrm K\) als Überschuss und \(\alpha/\Delta x^2=0{,}1\,\mathrm{s}^{-1}\) liefert die Modellfunktion \((1,-2,1)^\top\,\mathrm{K/s}\). Erst Euler etwa mit \(h=1\,\mathrm s\) erzeugt den nächsten Zustand \((1,8,1)^\top\,\mathrm K\).
\subt{Andere Raumdarstellungen führen zur selben Zeitfrage}
{\tblstart\begin{tabularx}{\linewidth}{@{}P{27mm}YY@{}}\toprule
\hd{Raumverfahren}&\hd{Was wird gespeichert?}&\hd{Woran prüfe ich das Ergebnis?}\\\midrule
Differenzen&Knotenwerte. Steigungen aus Nachbarn.&Fehlerabnahme bei Gitterverfeinerung.\\
Finite Volumen&Zellmittelwerte. Flüsse über gemeinsame Flächen.&Zu- und Abfluss jeder Zelle.\\
Finite Elemente&Koeffizienten lokaler Basisfunktionen in der schwachen Form.&Randbedingungen, Energie und Gitterfehler.\\
Spektralverfahren&Koeffizienten globaler Moden.&Abklingen hoher Moden bei glatten Daten.\\\bottomrule
\end{tabularx}\par}
Mit finiten Elementen entsteht oft \(M\dot U=r(t,U)\). Ist die Massenmatrix \(M\) invertierbar, ergibt das Lösen eines linearen Systems die Rate. Bleibt ein masseloser Rand als algebraische Unbekannte erhalten, gehören eine Nebenbedingung und passende Anfangswerte zum DAE-Problem.
\srcs{\href{https://fncbook.com/methodlines/}{Driscoll und Braun, Method of Lines} · \see{molrand}{Ein Rand mit algebraischer Bindung} · \see{schwach}{Schwache Form}.}

\lektion{Kann ein erhaltendes Verfahren negative Werte liefern?}{laxtransport}{Rechenversuch}{Transport}
Markierstoff bewegt sich in einem periodischen Kanal mit konstanter Geschwindigkeit \(a>0\). Vier gleich große Zellen enthalten anfangs die normierten Mittelwerte \((0,1,1,0)\). Während eines Schritts bewegt sich das Profil um eine halbe Zelle. Die dimensionslose Courant-Zahl ist damit \(C=ah/\Delta x=1/2\). Die linke Nachbarin von Zelle 1 ist Zelle 4.
\subt{Zwei Regeln, dieselben Daten}
Das Upwind-Verfahren nimmt den zufließenden Wert von links. Lax--Wendroff berücksichtigt zusätzlich die Krümmung aus beiden Nachbarn. Für die lineare Transportgleichung \(u_t+a u_x=0\) lauten beide vollständig diskreten Regeln
\begin{align*}
 U_j^{n+1}&=U_j^n-C(U_j^n-U_{j-1}^n),\\
 V_j^{n+1}&=V_j^n-\tfrac C2(V_{j+1}^n-V_{j-1}^n)
             +\tfrac{C^2}2(V_{j+1}^n-2V_j^n+V_{j-1}^n).
\end{align*}
{\tblstart\begin{tabularx}{\linewidth}{@{}Yrrrrr@{}}\toprule
\hd{Nach einem Schritt}&\hd{Zelle 1}&\hd{2}&\hd{3}&\hd{4}&\hd{Summe}\\\midrule
Upwind&0&0,5&1&0,5&2\\
Lax--Wendroff&$-0{,}125$&0,625&1,125&0,375&2\\\bottomrule
\end{tabularx}\par}
In Zelle 1 ist die Konzentration nach Lax--Wendroff negativ. Eine höhere Genauigkeitsordnung für glatte Daten beantwortet die Frage nach Positivität an einem Sprung also noch nicht. In diesem einen Schritt liefert Upwind sogar die exakten Zellmittel des um eine halbe Zelle verschobenen, ursprünglich stückweise konstanten Profils. Wiederholtes Mitteln verbreitert die Front.
\subt{Was können wir für beliebige Eingaben beweisen?}
\proofpart{Upwind bildet bei \(0\le C\le1\) die konvexe Kombination \((1-C)U_j+C U_{j-1}\). Sie liegt zwischen den Nachbarwerten. Bei periodischer Summation heben sich in beiden Regeln alle Differenzen auf. Beide erhalten daher die Gesamtmenge \(\Delta x\sum_j U_j\).}
\srcs{\href{https://www.clawpack.org/fvmhp_materials/}{LeVeque, Finite Volume Methods for Hyperbolic Problems} · \texttt{python methodenlabor.py} · \see{schluss6}{Fragen zur Methodenwahl}.}
